\section{Introduction}
\label{sec:introduction}

Rolling element bearings are fundamental components in modern industrial machinery, whose operational status directly determines the reliability and safety of the entire system \cite{lei2020deep_fault, jia2016sae_bearing, lei2020deep_fault}. Given that bearing failures can lead to catastrophic mechanical breakdowns, significant economic losses, and even human casualties, Condition Monitoring and Fault Diagnosis (CMFD) have become indispensable disciplines in the era of Industry 4.0 \cite{smith_cwru, zhang2021cwru_review, jia2016sae_bearing}. Traditional diagnostic approaches, relying on physical modeling and expert-defined signal processing features, have demonstrated effectiveness under stable laboratory conditions \cite{samanta2006svm_bearing, saidi2022triplet_svm, support2019svm_bearing, samanta2006svm_bearing}. However, with the increasing complexity of mechanical systems and the demand for automated monitoring, researchers have progressively shifted toward data-driven methodologies, particularly deep learning-based techniques, which can automatically extract high-level feature representations from massive volumes of raw vibration data \cite{chen2023cwru_deep, chen2021sae_fault, shao2018sae_bearing, chen2021sae_fault}. Consequently, the maintenance paradigm is evolving from reactive repairs to intelligent predictive maintenance, ensuring high availability of industrial assets across diverse operating scenarios \cite{lei2020deep_fault, chen2023cwru_deep, zhang2021cwru_review, lei2020deep_fault}. The continuous monitoring and automated analysis of these signals provide a crucial safeguard for the continuity of modern manufacturing processes \cite{smith_cwru, jia2016sae_bearing}.

Despite the rapid advancement of deep learning in fault diagnosis, a critical bottleneck remains: the data imbalance problem \cite{fan2021enhanced_gan, liu2024gan_imbalanced_review, qiao2021imbalanced_bearing}. In practical engineering environments, mechanical equipment typically operates in a healthy state for most of its service life, making normal operation data abundant \cite{qiao2021imbalanced_bearing, liu2024gan_imbalanced_review}. In contrast, fault samples are difficult and expensive to collect due to the risks associated with deliberate destructive testing and the stochastic nature of actual failures \cite{liu2020gan_smote_fault, zhao2022cnn_smote_bearing, liu2020gan_smote_fault}. This inherent scarcity of fault data relative to normal data leads to a highly imbalanced dataset distribution, where the majority class significantly outweighs the minority fault categories \cite{zhao2022cnn_smote_bearing, li2021wgan_gp_fault, qiao2021imbalanced_bearing}. Standard classification algorithms trained on such imbalanced datasets tend to be biased toward the majority class (normal), resulting in poor generalization and high misclassification rates for the critical minority classes (faults) which often carry the most vital information for safety \cite{qiao2021imbalanced_bearing, fan2021enhanced_gan, liu2024gan_imbalanced_review, fan2021enhanced_gan}. Therefore, developing robust data augmentation techniques to alleviate the data imbalance and enhance the diagnostic performance on rare fault categories is a paramount challenge for intelligent fault diagnosis systems \cite{liu2024gan_imbalanced_review, han2020wgangp_bearing, li2021wgan_gp_fault, han2020wgangp_bearing}. The ability to accurately identify these rare but critical fault states under varying load and speed conditions is essential for preventing unforeseen catastrophic events in complex industrial clusters \cite{qiao2021imbalanced_bearing, han2020wgangp_bearing}.

To address data imbalance, several traditional and generative solutions have been explored in recent years. Traditional oversampling methods, such as Synthetic Minority Over-sampling Technique (SMOTE) \cite{chawla2002smote} and Adaptive Synthetic Sampling (ADASYN) \cite{he2008adasyn}, generate new samples by interpolating between existing minority instances in the feature space \cite{zhao2022cnn_smote_bearing, chawla2002smote, he2008adasyn}. However, these methods often struggle to capture the complex, high-dimensional temporal dependencies and non-stationary characteristics inherent in vibration signals, potentially introducing noise or causing overfitting due to the simple linear interpolation mechanism \cite{liu2020gan_smote_fault, zhao2022cnn_smote_bearing, liu2020gan_smote_fault}. With the emergence of Generative Adversarial Networks (GANs) \cite{goodfellow2014gan}, researchers have applied variants like Deep Convolutional GAN (DCGAN) \cite{radford2015dcgan} and Wasserstein GAN (WGAN) \cite{arjovsky2017wgan} to synthesize artificial fault samples \cite{li2021wgan_gp_fault, li2022multi_mode_acgan, radford2015dcgan, arjovsky2017wgan}. While these GAN-based methods can learn the underlying distribution of fault data, they often require a certain amount of initial fault samples to begin the generation process and are prone to training instability or mode collapse \cite{gulrajani2017wgangp, meng2022iacgan, gulrajani2017wgangp, meng2022iacgan}. Furthermore, standard GAN architectures lack an explicit mechanism to transform knowledge from the informative normal domain to the sparse fault domain, limiting their effectiveness when the minority class is extremely scarce and making it difficult to maintain structural consistency between real and synthetic signals \cite{zhu2017cyclegan, alotaibi2021cyclegan_signal, zhu2017cyclegan, alotaibi2021cyclegan_signal}. As a result, many existing GAN-based diagnostic frameworks still struggle to produce high-fidelity vibration signals that can fool sophisticated expert-level classifiers while maintaining clear fault semantic signatures \cite{meng2022iacgan, li2022multi_mode_acgan, shao2018acgan_1dcnn}.

In this paper, we propose a novel data augmentation framework named Conditional Auxiliary Classifier Cycle-Consistent Generative Adversarial Network restrained by Wasserstein distance with Gradient Penalty (CAC-CycleGAN-WGP) to tackle the bearing fault diagnosis problem under extreme data imbalance. Unlike conventional GANs that generate samples from random noise, our approach leverages CycleGAN's dual-mapping mechanism \cite{zhu2017cyclegan, fernandez2022cyclegan_bearing, zhu2017cyclegan} to perform cross-domain translation from normal vibration signals to specific fault categories. We introduce a Conditional Auxiliary Classifier (CAC) \cite{odena2017acgan, meng2022iacgan, odena2017acgan, meng2022iacgan} to guide the generation process, ensuring that the synthesized signals not only retain the physical characteristics of the source domain but also strictly adhere to the target fault category labels. To enhance training stability and sample quality, we integrate the Wasserstein distance loss and Gradient Penalty (WGP) \cite{gulrajani2017wgangp, han2020wgangp_bearing}, which effectively mitigates the gradient explosion and vanishing problems common in cycle-consistent architectures \cite{zhang2022cwgan_gp, li2022cwgan_gp_pca, zhang2022cwgan_gp}. The inclusion of the WGP component allows the model to undergo more robust training with larger batch sizes and higher learning rates without compromising the fidelity of the generated mechanical signatures \cite{gulrajani2017wgangp, li2021wgan_gp_fault}. Furthermore, we address the evaluation challenge by implementing a Stacked Autoencoder (SAE) \cite{hinton2006sae, shao2018sae_bearing, jia2016sae_bearing} for feature quality assessment and a Support Vector Machine (SVM) \cite{cortes1995svm, support2019svm_bearing, samanta2006svm_bearing} for final fault classification across two benchmark bearing datasets. This multi-stage framework ensures that the augmented data is not only statistically similar to the real data but also diagnostically valuable for downstream identification tasks \cite{shao2018sae_bearing, support2019svm_bearing}.

The main contributions of this work are summarized as follows:
\begin{enumerate}
    \item We propose the CAC-CycleGAN-WGP framework, which is the first to integrate conditional auxiliary classification into a cycle-consistent architecture for the cross-domain translation and augmentation of bearing vibration signals from normal to multiple fault domains.
    \item We incorporate the Wasserstein distance with gradient penalty into the CycleGAN training objective, significantly improving the stability of the dual-domain mapping and ensuring the generation of high-fidelity fault signatures even with limited training iterations.
    \item We establish a hierarchical evaluation scheme using SAE for high-dimensional feature reconstruction analysis and SVM for diagnostic performance verification, providing a comprehensive assessment of the augmented data's quality beyond simple visual inspection.
    \item Extensive experiments on the Case Western Reserve University (CWRU) and another dataset demonstrate that our method consistently outperforms state-of-the-art GAN variants and traditional oversampling techniques across various imbalanced ratios, specifically in low-sample regimes.
\end{enumerate}

The remainder of this paper is organized as follows. Section II presents the theoretical background of GAN, CycleGAN, and Wasserstein distance. Section III describes the proposed CAC-CycleGAN-WGP architecture and the detailed data augmentation workflow. Experimental results, sensitivity analysis, and comparative studies are provided in Section IV. Finally, Section V concludes the paper and discusses future research directions for intelligent fault diagnosis.
